{\large\textbf{E – Physics of piezoelectricity (20 pts)}}

Some crystalline, electrically insulating materials, such as quartz and lead
zirconate titanate, exhibit an electrical response when mechanical pressure is
applied. In short, mechanical stress polarizes the crystals, which is called
the piezoelectric effect. This phenomenon can be explained by the special
structure of their molecules: the deformation gives each molecule an electric
dipole moment. Conversely, mechanical stress is generated in the presence of an
electric field, a phenomenon known as the reverse piezoelectric effect.
However, in this problem, we will neglect the reverse piezoelectric effect.

This problem investigates a simple device that relies on piezoelectricity, the
piezoelectric element. It consists of the piezoelectric material placed between
two circular metallic plates. When a force, perpendicular to the metallic
plates, is applied to the element, a force-dependent voltage appears between
the plates.

\textbf{Equipment (see also Fig. 1)}

A\quad Piezo element with electrical leads, attached to a wooden base plate and
a cover plate with a small hole. The hole goes through the entire cover plate,
such that a small part of the top electrode of the piezo element is visible
through it. This electrode is thin and flexible.

B\quad Multimeter (inner resistances are given, refer to the important remarks
on the next page).

C\quad $1.5\,\mathrm{V}$ AA battery with connector.

D\quad Capacitor of unknown capacitance, with a diode soldered to one leg (when
a forward current flows through the diode, the voltage drop across it is
$0.56\,\mathrm{V}$).

E\quad 4 electric push-button switches (connected while pressed) on wires.

F\quad 6 crocodile clamps.

G\quad Digital scales (up to $10\,\mathrm{kg}$).

H\quad Digital stop watch.

I\quad 2 different rubber balls.

J\quad Wooden stand with adjustable release mechanism.

K\quad $50\,\mathrm{cm}$ ruler.

L\quad Small, wooden sticks (diameter of $2\,\mathrm{mm}$), in two different
lengths.

M\quad Large metal screw.

N\quad Wooden clothes peg.

O\quad Pencil, pen and pencil sharpener.

\textbf{Task E.1 — Elasticity of the ball (2.0 pts)}

Of the two rubber balls you are provided with, one is more elastic than the
other.

For the more elastic rubber ball, determine the fraction of kinetic energy
which is lost during a collision with a solid surface. Determine this fraction
for three different values of the initial kinetic energy.

\textbf{Task E.2 - Piezoelectric properties (10.0 pts)}

a\quad Measure the capacitance $C$ of the capacitor (Fig. 3b). (2.0 pts)

b\quad The metallic plates on the sides of the piezo element also act as a
capacitor. Find the capacitance $C_p$ of the piezo element. (2.5 pts)

c\quad Measure and plot how the voltage between the plates of the piezo element
depends on the total perpendicular force, which is evenly distributed over the
surface of the piezo via its wooden cover plate. For low forces, the dependence
is linear; find the slope $\beta$ in this regime. (4.0 pts)

d\quad The molecules of the crystals can only have polarizations lower than a
certain critical value. Find the maximal (saturation) voltage of the piezo, the
pressure $p_{sat}$ at saturation and the maximal surface density
$\sigma_{max}$ of the charge on the surface of the piezo element. (1.5 pts)

\textbf{Task E.3 - Small area behaviour (1.0 pts)}

When a force is applied to a small region of the piezo crystal, due to
electro-mechanical coupling, the crystal will try to curve. However, the wooden
plates will prevent this, and as a result, mechanical stress will appear in
other parts of the crystal, too.

How much will the electrical response change when the force is applied to a
small area of the crystal? Consider only the linear range of the response.

\textbf{Task E.4 - Deformation of the ball (4.5 pts)}

In this part you will be dropping the more elastic rubber ball on the piezo
element. During the collision between the ball and the piezo, the ball
experiences deformations. You can assume that the force $F$ acting on the ball
depends on the ball's elastic deformation $x$ as a power law:
\begin{equation}
F = kx^\alpha \ . \tag{1}
\end{equation}

Find the exponent $\alpha$ and the material constant $k$.

\textbf{Task E.5 - Interaction time (2.5 pts)}

With the result from the previous task, it would be possible to determine the
interaction time $\tau$ of the more elastic ball with the wooden surface.
However, for the less elastic ball, there is no simple description such as
Eq. 1. Alternatively, we can make the following assumption.

If, for a certain collision speed $v_0$, the force experienced by the ball as a
function of time can be described as $F_0(t) = f(t)$, then for any other speed
$v_1$, the time dependence takes a similar shape and we can express the force
as:
\begin{equation}
F_1(t) = a_1 f(b_1 t) \ . \tag{2}
\end{equation}

Estimate and plot how the interaction time $\tau$ scales with the collision
speed $v$ for collisions of the less elastic rubber ball with a solid surface.

\textbf{Important practical remarks}

\begin{itemize}
\item It is expected that you provide circuit diagrams with all the electrical
measurements you take. Use the symbols provided in Fig. 2.
\item Be careful not to short the multimeter! The internal resistances given
are only valid for DC voltage measurements.
\item Do not exceed a total load of $100\,\mathrm{N}$ of force on the piezo
element.
\item To connect two wires, it is recommended to wrap them around each other
and use the crocodile clamps to secure the connection, see Fig. 3a.
\item The diode's polarity is visible in Fig. 3b.
\item You can adjust the height of the release mechanism via the two screws on
the back, see Fig. 4. You can also adjust the horizontal position by
unscrewing the vertical nut at the release mechanism.
\item Be careful not to let the bouncy ball escape. You can minimize the chances
of the ball leaving your table by placing the stand with the release mechanism
against one of the walls or in a corner of your desk and use the ruler as
another, transparent wall that is secured with the clothes peg. If you still
loose the ball, ask for assistance, \textbf{do not leave your place}.
\item After a certain time of inactivity, the multimeter will start to beep
and, after another few minutes, will shut down. You can prevent shut-down by
pressing any button.
\item Be aware that some regions of your table may spring back under load. It
is wise to use stable regions of your table for the bouncing experiment.
\item For the evaluation of a series of measurements, you are expected to plot
the data.
\item Only those measurements and evaluation methods that promise the highest
precision and accuracy are awarded full marks. Choose your approaches
accordingly while keeping in mind the precision of your tools. However, no
marks are awarded for error estimation.
\item When measuring DC voltages with the multimeter in a range with the
smallest digit resolution being $\delta U$, the uncertainty $\Delta U$ of the
measured value $U$ can be estimated via:
\begin{equation}
\Delta U = 0.7\% \cdot U + 3 \cdot \delta U \ . \tag{3}
\end{equation}
\item The internal resistance of your multimeter in the DC voltmeter setting
depends on the measurement range and model number. In Fig. 5, you can see
where to find the model number of your voltmeter. The following table provides
the internal resistances for two measurement ranges and all models.
\end{itemize}

\begin{center}
\begin{tabular}{ccc}
\textbf{model number(s)} & \textbf{range} & \textbf{internal resistance} \\
all models & $2\,\mathrm{V}$ & $11.1\,\mathrm{M\Omega}$ \\
UT33B+ & $200\,\mathrm{mV}$ & $9.65\,\mathrm{M\Omega}$ \\
UT33C+ & $200\,\mathrm{mV}$ & $9.91\,\mathrm{M\Omega}$ \\
UT33D+ & $200\,\mathrm{mV}$ & $9.70\,\mathrm{M\Omega}$ \\
\end{tabular}
\end{center}

\begin{center}\includegraphics[width=0.53\linewidth]{figures/EuPhO_2024_Q4_fig1.jpg}\end{center}

Figure 1: Overview of all the materials used in this experiment.

\begin{center}\includegraphics[width=0.55\linewidth]{figures/EuPhO_2024_Q4_fig2.png}\end{center}

Figure 2: The most common circuit symbols: \textbf{A} - resistor, \textbf{B} -
ammeter, \textbf{C} - voltmeter, \textbf{D} - ohmmeter, \textbf{E} - switch,
\textbf{F} - battery, \textbf{G} - capacitor, \textbf{H} - diode, \textbf{I} -
piezo element. Wires marked with a red star have a red coating.

\begin{center}
\begin{minipage}[t]{0.30\linewidth}\raggedright
\includegraphics[width=\linewidth]{figures/EuPhO_2024_Q4_fig3.jpg}

(a) Recommended electrical connection with wires wrapped around another and
secured with a crocodile clamp.
\end{minipage}\hspace{0.01\linewidth}%
\begin{minipage}[t]{0.23\linewidth}\raggedright
\includegraphics[width=\linewidth]{figures/EuPhO_2024_Q4_fig4.jpg}

(b) The diode attached to the capacitor with the circuit symbol placed next to
the diode in correct orientation.
\end{minipage}
\end{center}

Figure 3: Helpful remarks on the electric components

\begin{center}\includegraphics[width=0.53\linewidth]{figures/EuPhO_2024_Q4_fig5.jpg}\end{center}

Figure 4: The two wingnuts on the back can easily be loosened to allow the
release mechanism to slide vertically. Also, the vertical nut can be used to
enable horizontal adjustment.

\begin{center}\includegraphics[width=0.47\linewidth]{figures/EuPhO_2024_Q4_fig6.jpg}\end{center}

Figure 5: The red arrow points towards the model number of the multimeter.
